\documentclass[a4paper,11pt]{article}
\usepackage{graphicx,geometry,subcaption}
\usepackage{fontspec}
\setmainfont{Times New Roman}


\usepackage{hyperref}
\hypersetup{
	colorlinks=true,
	urlcolor=blue,
	linkcolor=black,
	citecolor=black
}

\usepackage{setspace}
\singlespacing

\usepackage[spanish]{babel}
\geometry{margin=2.5cm}

\usepackage{fancyhdr}
\pagestyle{fancy}

\fancyhf{}
\fancyfoot[C]{\thepage}
\renewcommand{\headrulewidth}{0pt}

\usepackage{titlesec}
\titleformat{\section}
	{\normalfont\bfseries\fontsize{11}{13}\selectfont}
	{\thesection}
	{1em}
	{\MakeUppercase}
	
\titlespacing*{\section}
	{0pt}
	{11pt}
	{11pt}

\titleformat{\subsection}
	{\normalfont\bfseries\fontsize{11}{13}\selectfont}
	{\thesubsection}
	{1em}
	{}

\titlespacing*{\subsection}
	{0pt}
	{11pt}
	{11pt}


\newcommand{\practicatitle}[9]{
    \begin{center}
        \vspace*{11pt}
        
        {\bfseries\MakeUppercase{#1}\par}
        
        \vspace{18pt}
        
        {\bfseries #2\par}
        
        \vspace{6pt}
        
        #3\par
        #4\par
        #5\par
        #6\par
        
        \vspace{12pt}
        
        #7\par
        #8\par
        #9\par
        
        \vspace{18pt}
    \end{center}
}

\begin{document}
\practicatitle
{Diseño y Desarrollo de un gateway Ethernet-5G privado}
{Williams, Robbie Keith – X8503344Z (Estudiante)}
{García Sánchez, Antonio Javier (Tutor de Empresa)}
{e-mail (tutor de empresa): \href{mailto:antoniojavier.garcia@upct.es}{antoniojavier.garcia@upct.es}}
{Universidad Politécnica de Cartagena}
{Antiguo Cuartel de Antigones. Plaza del Hospital, 1. 30202 Cartagena}
{García Sánchez, Antonio Javier (Tutor Académico)}
{e-mail (tutor académico): \href{mailto:antoniojavier.garcia@upct.es}{antoniojavier.garcia@upct.es}}
{Universidad Politécnica de Cartagena, Departamento de Tecnologías de la Información y las Comunicaciones}
\tableofcontents
\newpage
\section*{Resumen}
En esta práctica, me he formado en programación de sistemas embebidos, específicamente programación en ESP32 utilizando el entorno ESP-IDF. El objetivo del proyecto en el cual estoy trabajando es el diseño y desarrollo de un gateway Ethernet-DECT que permita recibir comunicación por ethernet y retransmitir los datos recibidos por el protocolo dect y vice-versa.También he aprendido a estudiar diseños de PCBs a través de herramientas como KiCad y FreeCAD.
\addcontentsline{toc}{section}{RESUMEN}
\section{Descripción concreta y detallada de las tareas, trabajos desarrollados y departamentos de la entidad a los que ha estado asignado.}
Las tareas que he desarrollado y estoy desarrollando a lo largo de este puesto son:
\begin{itemize}
	\item Entender el entorno ESP32, leyendo la documentación y estudiando las guías de referencia API además de preguntarle dudas a mis compañeros.
	\item Observar y analizar el diseño de la placa en KiCad y FreeCAD, con el objetivo de mejorar mi entendimiento del proyecto y de sus requisitos y objetivos. Colaborar con compañeros de hardware para identificar posibles problemas y encontrar soluciones.
	\item Estudiar el código del firmware de la placa para entender su estructura y las funciones de cada componente.
	\item Hacer el build del proyecto y flashearlo en una placa de prueba para comprobar su funcionalidad.
	\item Programar en C un firmware simulado de la placa para poder comprobar su correcta estructura y depurar posibles errores, flasheando y analizando los resultados de su ejecución en la placa de prueba.
\end{itemize}
Todas estas tareas se realizaron en el laboratorio de software Qartia, situado en la primera planta del edificio de Antigones.
\section{Valoración de las tareas desarrolladas con los conocimientos y competencias adquiridos en relación con los estudios universitarios.}
Las competencias adquiridas a lo largo del grado en Ingeniería Telemática que he aplicado en las prácticas son principalmente dos:
\begin{enumerate}
	\item \textbf{Programación en C: } Aunque a lo largo de la carrera de Ingeniería Telemática no tenemos ninguna asignatura que aborde C o C++ mas allá de lo más básico, los conceptos básicos aprendidos en Java se pueden aplicar también en la programación del firmware del ESP32 en C.
	\item \textbf{Entorno ESP32/ESP-IDF: } En el curso de mis estudios nunca he trabajado con ESP32, pero he podido aplicar los conocimientos aprendidos en la asignatura de Sistemas Digitales Basados en Microprocesadores sobre el proceso de build, flash y posterior depuración de la ejecución del programa sobre la placa física.   
\end{enumerate}
\section{Relación de los problemas planteados y el procedimiento seguido para su resolución.}
\begin{itemize}
	\item \textbf{Desarrollo del firmware: } Para el desarrollo de la práctica es importante la implementación y validación del firmware del proyecto. El procedimiento seguido fue el siguiente:
	\begin{enumerate}
		\item \textbf{Estudio del entorno ESP-IDF: } Para poder estudiar bien el código heredado en el proyecto, primero fue importante entender el propio entorno ESP-IDF. Sus APIs, métodos y funciones básicos, el manejo de errores, etc. Aunque en ESP-IDF se utiliza C, como todos los proyectos ESP dependen de esta librería, la manera de programar y la estructura del código son diferentes, por lo tanto es importante entender al menos lo más básico para poder trabajar con este framework.
		\item \textbf{Estudio del código: } En segundo lugar, realicé un estudio del código heredado, para entender la funcionalidad de cada componente y el comportamiento del firmware. 
		\item \textbf{Pruebas en placa física: } El siguiente paso que tomé fue buscar una placa ESP32 en la cual realizar pruebas y hacer modificaciones al firmware. Primero probé varios ejemplos incluidos en el propio entorno ESP-IDF para comprobar el correcto funcionamiento de la placa obtenida gracias a los compañeros del laboratorio de hardware. Cuando comprobé que los ejemplos se ejecutaban de forma correcta y la salida en consola era la esperada, continué con el siguiente paso.
		\item \textbf{Prueba del firmware en placa de desarrollo: } Tras realizar las pruebas con los proyectos de ejemplo aportados por el entorno ESP-IDF, decidí probar el firmware del gateway ethernet-DECT en la placa de desarrollo. Aquí fue donde pude encontrar el primer gran problema del proyecto. El firmware fallaba a la hora de inicializar el link ethernet. Consulté con mis compañeros, sobretodo con Julián, la forma más eficiente de depurar el firmware, ya que yo estaba acostumbrado a utilizar debuggers integrados como el de \href{https://code.visualstudio.com/docs/debugtest/debugging}{Visual Studio Code} (o parecidos) y herramientas como \href{https://sourceware.org/gdb/}{GDB}, los cuales no se pueden utilizar en el entorno de desarrollo ESP32. Gracias a la aportación de mis compañeros e investigación propia, entendí que la forma más viable de depurar el firmware en este contexto es simplemente con sentencias \href{https://docs.espressif.com/projects/esp-idf/en/stable/esp32/api-reference/system/log.html#_CPPv47esp_log16esp_log_config_tPKcPKcz}{ESP\_LOG}, que imprimen en pantalla el valor obtenido de una función o una expresión además del valor esperado para así comprobar la causa de los errores indicados en el build. Tras varias horas de investigación y depuración, encontré la causa del problema. El proyecto que heredé asumía cierto chip de ethernet, el IP101GRI, mientras que la placa de desarrollo que tenía disponible consta de un chip distinto, el W5500. El código de cada chip es incompatible con el otro, y hacer una migración del código entero a otro controlador de ethernet no parecía viable en su momento, entonces decidí tomar otro camino.
		\item \textbf{Desarrollo de firmware simulado: } En vez de intentar ejecutar un firmware diseñado para otra placa, con características distintas y chips/drivers diferentes, decidí simular el comportamiento del firmware para comprobar que la lógica del código es correcta. Para esto, elaboré un código "mock" para cada archivo de código del firmware. Es decir, un código que simula el comportamiento del hardware físico. Este código simplemente asume que el hardware funciona correctamente y devuelve valores predefinidos para simular el comportamiento real. Así pude depurar la lógica del programa y encontrar fallos en su planteamiento y/o ejecución. El proceso de desarrollo empezó con un estudio más profundo del código, tomando nota de el objetivo de cada función y variable global del código. A partir de esta información empecé desde cero a escribir el firmware simulado, reemplazando cada función requerida por una versión "MOCK" simulada, empezando con el 'app\_main()'. En cuanto el 'mock\_app\_main()' fue completado con todas las llamadas que establece, repetí el mismo proceso con cada componente del firmware. Esto es:
		\begin{itemize}
			\item gw\_bridge
			\item gw\_dect
			\item gw\_diag
			\item gw\_eth
			\item gw\_watchdog
		\end{itemize}
		En cuanto la elaboración del firmware simulado para el resto de componentes se completó, empecé un largo periodo de depuración.
		\item \textbf{Depuración del firmware simulado: } inicialmente el build con el firmware simulado daba muchos errores de compilación. Aunque encontré y arreglé algunos errores sintácticos en el código desarrollado, no conseguí completar el build y no entendía la causa. Tras revisar de nuevo la documentación y la referencia API de espressif, encontré que una posible causa sería tener que modificar los archivos CMakeList.txt dentro de cada carpeta del firmware. Tras adaptar estos archivos para tener en cuenta el firmware simulado en vez de el original, pude completar el build e intentar flashearlo en la placa de desarrollo.
		\item \textbf{Validación del firmware simulado en la placa de desarrollo: } El flash del firmware simulado funcionó correctamente con el primer intento y conseguí verificar la lógica del firmware con la salida en consola. Por culpa de no implementar pequeños retardos en el firmware simulado para considerar los retardos reales que pueda tener realizar varias funciones como asignar una IP al gateway Ethernet, inicializar el driver DECT, entre otros, el programa se ejecuta muy rápidamente y cuesta ver el resultado en pantalla. Hice un pipe en bash de la salida a un archivo .txt y en este archivo pude verificar que la lógica del firmware es correcto.
	\end{enumerate}
	\item \textbf{Estudio de la placa diseñada en KiCad: } En paralelo con la validación y desarrollo del firmware, realicé un estudio de la placa diseñada a través de las herramientas 2D y 3D del entorno KiCad, con el fin de entender el comportamiento del diseño. Con la ayuda de mis compañeros del laboratorio de hardware, pude identificar varios problemas con el diseño que tendrían que ser rectificados posteriormente. Un chip que atravesaba el PCB, elementos desconectados, vías localizadas en los pads, entre otros. 
\end{itemize}
\section{Anexos.}
\begin{itemize}
	\item \href{https://docs.espressif.com/projects/esp-idf/en/stable/esp32/api-reference/index.html}{Documentación espressif}, utilizada como referencia principal para comprobar APIs de desarrollo, manejo de errores, etc.
	\item \href{https://docs.espressif.com/projects/esp-dev-kits/en/latest/esp32/esp32-ethernet-kit/user_guide.html}{ESP32-Ethernet-Kit v1.2}, placa usada para flashear el firmware y realizar las pruebas necesarias.
	\item \href{https://github.com/Robinhood-700/gateway-eth-5g-privado}{Repositorio GitHub} en el cual se almacena el código fuente del proyecto además de los resultados obtenidos.
	\item \href{https://github.com/espressif/esptool}{ESPtool}, software utilizado para realizar el build del proyecto y realizar el proceso de flash.
\end{itemize}
\thispagestyle{fancy}
\fancyfoot[L]{PDF Compilado con \LaTeX}
\fancyfoot[C]{\thepage}
\fancyfoot[R]{\today}
\end{document}
